\MajorSection{القسم الثاني}

\Couplet{١٨}{في هذه القفار الكئيبة من أرض ولدت من البحر، في هذه البراري التي لا يسكنها إلا هو؛}{أي ماضٍ شبحي يعود إلى الحياة، وأي تعاقب للسنين نبصر!}{}

\Couplet{١٩}{نحدّق وراء الخط الأزرق الرقيق الذي يطوّق دائرة الأفق البعيد؛}{فلماذا تطارد بصرنا المحزون هذه الأشباح، ومن أين تنهض هذه الظلال الطيفية؟}{}

\Couplet{٢٠}{أي أسئلة لا تنتهي تقلق الفكر: من أين، وإلى أين، ومتى، وكيف؟}{وأي سعي مولع أحمق إلى قراءة الكتاب المنقوش على جبين الإنسان!}{}

\Couplet{٢١}{وها نحن نقف على رأس نقطة من الزمان، بين أزليّة وأبديّة؛}{وأسرارهما المروّعة تلتف حولنا، وتثقل أعيننا بسوادها العميق.}{}

\Couplet{٢٢}{«هذا الليل المظلم، وهذه الأمواج المرعبة، وهذه الرياح والدوامات الصاخبة المخيفة؛}{ما بال الذين يخطون مطمئنين على شاطئ الأمان يكترثون لبؤس حالنا؟»}{}

\Couplet{٢٣}{هكذا قال شاعر الحب والخمر، الذي لم يرتفع حلمه بالجنة}{فوق كأس الكوثر المترعة، والحور ذوات الأعين الشديدة البياض والسواد.}{\BurtonNote{حاشية برتون: حافظ الشيرازي.}}

\Couplet{٢٤}{وا أسفاه! إن شوط عمري، وهو ستون سنة، قصير؛ لكنه طويل بما يكفي}{ليُسأم إحساسي من مسرّات كهذه لا سرور فيها، من الحب والحور والخمر وكل ما معها.}{}

\Couplet{٢٥}{ويتفاخر آخر بأنه سيطلّق العقل العجوز العقيم من فراشه؛}{ويتزوج بنت الكرمة مكانه؛ حمقى هم الذين يصدقون كلمة مما قال!}{\BurtonNote{حاشية برتون: عمر الخيام، شاعر فارس وصانع الخيام.}}

\Couplet{٢٦}{ويصيح الصوفي: «قولُهم إنك تراب وإلى التراب تعود لم يُقَل قط في الروح الإنسانية»؛}{طوبى لمن يملك مثل هذه الموهبة، فيسأل عن غايتها.}{}

\Couplet{٢٧}{«أهذا هو كل شيء؟ ألهذا نولد، لنبكي قليلًا ثم نموت؟»}{هكذا يغني الشاعر الضحل الذي لا تزال حياته تكدح عند حرف «أنا».}{}

\Couplet{٢٨}{«لم تسمع أذن، ولم تر عين، نعيم الذين يدخلون ملكوتي السماوي»؛}{هكذا قال عيسى، الذي ناح على أحزاننا وخطيئتنا.}{}

\Couplet{٢٩}{كلمات أكثر مما ينبغي، أو أقل مما ينبغي! وأي شيء أيسر على ألوهيتك}{من لمحة صغيرة من الفردوس تفتح بها عيني الإنسان وأذنيه؟}{}

\Couplet{٣٠}{«أنا الحق! أنا الحق!» نسمع العارف السكران بالله يصيح؛}{«العالم الصغير قائم فيَّ؛ والله الأزلي ليس إلا أنا!»}{}
{\EditorNote{حجي عبدو اليزدي يقتبس قولاً غير مألوفٍ هنا وينسبه للحلاج. المترجم بورتون رغم درايته ببعض ما نسب إلى المتصوفة فهو يعطي مثلاً "وليس في جبّتي إلا الله" وينسبه للحلاج.  لكن الداري يعلم أن هذي إحدى الشطحات المنسوبة إلى أبي يزيد البسطامي، والأدرى يظن أنها منسوبةٌ أيضاً  إلى غيره. من أقرب الأقوال حرفاً القول المنسوب للإمام علي "وَتَحْسَبُ أَنَّكَ جِرْمٌ صَغِيرٌ وَفِيكَ انْطَوَى الْعَالَمُ الأَكْبَر". أحد أقرب الأقوال المنسوبة للحلاج حكماً "أنا مَن أهوى ومَن أهوى أنا نحن روحانِ حَلَلْـنا بدنا."}}
%كامل مصطفى الشيبي في شرح ديوان الحلاج ص. ٨٩ ينسبها أيضاً لأبي سعيد بن أبي الخير.

\Couplet{٣١}{كان منصور حكيمًا؛ لكن الذين ضربوه بالحجارة المقذوفة كانوا أحكم؛}{ومع أن دمه أدّى الشهادة، لم تستطع قوة الحكمة أن تصلح عظامه.}{\BurtonNote{حاشية برتون: متصوف شهير رُجم بسبب التجديف.}}

\Couplet{٣٢}{«كُل واشرب والْهُ؛ فما تبقّى من الحياة لا يساوي نقرة إصبع»؛ هكذا قال الملك؛}{وأحسب أن هذا القول يقول أكثر مما ينبغي؛ فالخنزير سيقول الشيء نفسه!}{}

\Couplet{٣٣}{بهائم على قدمين ترعى في الحياة، وقد أُعدّت بالموت لتصير ترابًا؛}{تنحني إلى الأرض التي منها كانت، وتجد هناك لذّاتها اللائقة بها.}{}

\Couplet{٣٤}{أما أنتم، من مادة أرق وأشرف، يا من يقودكم العالي إلى ما هو أعلى؛}{فأي رابطة مشتركة تربط قلوبكم بمخلوقات الزريبة وحظيرة الخنازير؟}{}

\Couplet{٣٥}{«في رجاء يقيني بالحياة الآتية أجتاز هذا المشهد المتبدّل»؛}{يزمجر الزاهد، ويتهادى في وادي دموعه بطلعة واثقة.}{\BurtonNote{حاشية برتون: صاحب الاعتقاد «المحترم» من طراز الـ \textenglish{Philister}، أي ضيق الأفق المتشبث بالمألوف.}}

\Couplet{٣٦}{أأنت أحكم من ابن عمران، وأنت تعرف عالم الغد هذه المعرفة؛}{وتعرف المستقبل، فيما الماضي غير موجود، والحاضر مجرد حلم؟}{\BurtonNote{حاشية برتون: موسى كما يرد في القرآن.}}

\Couplet{٣٧}{ماذا تعرف، أيها الإنسان، عن الحياة؟ ومع ذلك، وأنت دائمًا بين الرحم والقبر؛}{تثرثر عن الحياة المقبلة، ولا بد أن تهذي بالجنة والنار.}{}

\Couplet{٣٨}{العالم قديم وأنت فتي؛ العالم كبير وأنت صغير؛}{فكفّ، يا ذرة عمرها لحظة، عن أن تعدّ نفسك الكل في الكل!}{}
